% gerado por ferramentas/md2tex_enunciado.py a partir de questoes/q104.md (não editar à mão)
A água é uma substância essencial à vida e pode ser encontrada na natureza em três estados físicos: sólido, líquido e gasoso. As propriedades físico-químicas que definem cada um desses estados estão representadas no diagrama de fases, em que as curvas em vermelho representam o estado de equilíbrio entre duas fases. Por meio dele, é possível determinar qual é o estado físico de uma substância em um dado valor de pressão (p) e de temperatura (T).\par
\figq[0.27\linewidth]{q104_fig1.png}
\fonte{ATKINS, Peter; JONES, Loretta. \textbf{Princípios de química:} questionando a vida moderna e o meio ambiente. Porto Alegre: Bookman, 2006 (adaptado).}
Qual é a mudança de estado físico que ocorre com a água quando, à pressão constante, a temperatura varia de A para B ao longo do trajeto tracejado?\par
\begin{alternativas}
  \item Fusão, a água muda do estado sólido para o líquido.
  \item Ebulição, a água muda do estado líquido para o gasoso.
  \item Sublimação, a água muda do estado sólido para o gasoso.
  \item Solidificação, a água muda do estado líquido para o sólido.
  \item Condensação, a água muda do estado gasoso para o líquido.
\end{alternativas}
